\documentclass[letterpaper,10pt]{texMemo}

\usepackage{graphicx}

\memoto{Dr. Mark Yoder}
\memofrom{Michael McDonald}
\memosubject{Week 4 (GDB, Make) Memo}
\memodate{\today}
\logo{\includegraphics[width=0.5\textwidth]{RoseLogo.png}}

\begin{document}
\maketitle

\section{Exercise Summary}

\noindent Matt Skorina and I have Homework 4 almost complete (we have all the
code written for searching and following, but the stepper is not stepping
properly).

\begin{center}
\begin{tabular}{|p{4cm}|l|l|p{8cm}|}
\hline
\textbf{Objective} & \textbf{Exercise} &\textbf{Status} & \textbf{Notes}\\
\hline
GDB & Exercise 14 & Completed & GDB is powerful but can be difficult to work
with, as commands aren't always intuitive. \\
\hline
Make & Exercise 15 & Completed & Makefiles are awesome. I feel like they can
get complicated quickly for large projects, but that is still way simpler than
trying to do it all by hand.\\
\hline
\end{tabular}
\end{center}

\noindent The main points I took away from these exercises were the following:
\begin{itemize}
 \item There are some pretty awesome software tools out there (that are open
source) that allow you to simplify your workflow by creating a solid tool chain.
 \item That being said, creating your tool chain can be somewhat difficult and
also time consuming. Luckily, these tools were already installed and ready to
use.
 \item LLDB, the ``new'' version of GDB, is actually much nicer. I used it with
Objective-C programming this summer, and I found it much nicer to use and just
as powerful (if not more so).
\end{itemize}

\end{document}
